\begin{abstract}
Nominal payload ratings for articulated robots are typically derived from worst-case configurations, resulting in uniform payload constraints across the entire workspace. This conservative approach severely underutilizes the robot's inherent capabilities---our analysis demonstrates that manipulators can safely handle payloads well above nominal capacity across broad regions of their workspace while staying within joint angle, velocity, acceleration, and torque limits. To address this gap between assumed and actual capability, we propose a novel trajectory generation approach using denoising diffusion models that explicitly incorporates payload constraints into the planning process. Unlike traditional sampling-based methods that rely on inefficient trial-and-error, optimization-based methods that are prohibitively slow, or kinodynamic planners that struggle with problem dimensionality, our approach generates dynamically feasible joint-space trajectories in constant time that can be directly executed on physical hardware without post-processing. Experimental validation on a 7 DoF Franka Emika Panda robot demonstrates that up to $67.6\%$ of the workspace remains accessible even with payloads exceeding $3$ times the nominal capacity. This expanded operational envelope highlights the importance of a more nuanced consideration of payload dynamics in motion planning algorithms.
\end{abstract}